\section{Discussion}
\label{sec:discussion}

\paragraph{What carries the performance}
The ablations separate the two kinds of structure that \gs{} combines. The chemistry-aware \emph{embedding}, which the snap operator searches, carries essentially all of \gs{}'s advantage over random search. That advantage appears only on O2, and replacing the embedding with geometric descriptors removes it. The chemistry-aware \emph{communication topology} adds little on top. Randomising it, replacing it by a Watts--Strogatz graph, or switching off the neighbour term changes recall by at most a few thousandths, with small and inconsistent effect sizes. A plausible reason for the objective dependence is that the lowest PBE gaps are dominated by open-shell transition metals (Cu, Co and Fe make up 41 of the 87 O2 hits), which the metal block of the embedding separates well. The O1, O3 and O4 hits are spread over 43--61 Leiden communities with no comparable descriptor signature.

\paragraph{Homophily is necessary but not sufficient}
Gate~4 showed strong band-gap homophily on MOFGalaxyNet ($r\approx0.5$--0.66). Homophily alone does not make a graph a good communication structure for a population optimiser, for three reasons. First, much of it comes from building-block duplicates that sit in near-cliques. Two agents in the same clique snap to nearly identical MOFs in the embedding, so the neighbour term carries little information beyond what the snap operator already uses. Second, \gs{} agents use the graph only to decide whom to listen to; the positions they average are embedding coordinates, and similar MOFs are already close in that space. Third, the graph is fragmented (1{,}075 isolated MOFs and a 42\% giant component at $\phi^\ast$), so agents sitting on isolated MOFs or small components have no neighbours to learn from. The greedy walk, which follows edges directly, has the lowest recall of the adaptive methods on O1 at the 2\% and 5\% budgets, which suggests that graph locality is a weak guide on that objective.

\paragraph{Position relative to baselines}
Surrogate-based acquisition, especially the bootstrap random-forest Thompson sampler, dominates every objective and budget, as is typical in finite-pool materials search. The snap-based evolutionary methods (PSO, DE, GA) also outperform \gs{}. One possible reason, which we did not test directly, is that SOCIAL's synchronisation and elite terms pull agents towards a common region, so successive snaps fill the neighbourhood of the elite. The ablations that switch these terms off did not change recall significantly, however, so this explanation remains tentative. As an acquisition policy, \gs{} is therefore cheap and better than random on objectives aligned with the embedding, but not competitive with surrogate models.

\paragraph{Reproducibility of MOFGalaxyNet}
The Gate~1 audit has implications beyond this study. The published network statistics correspond to a threshold near 0.7 rather than the stated 0.9, and the released code uses a different fingerprint and weighting from those described in the paper. We recommend that studies building on MOFGalaxyNet use the recipe and threshold reproduced here, or report both.

\paragraph{Limitations}
(i) The oracle is a database look-up. Results transfer to real DFT campaigns only to the extent that QMOF's single-linker, single-metal subset (8{,}697 of 20{,}372 entries) is representative. (ii) Topology was varied while the 32-d embedding and all SOCIAL defaults were held fixed. A different coupling, for example moving agents along graph edges instead of in an embedding, might use the topology more directly. (iii) The ensemble Thompson sampler uses random forests rather than a GNN, which may understate what model-based acquisition can achieve. (iv) The pilot had 10 seeds and about 29\% power for its observed effect, and the 30-seed comparison still has 20--64\% power. Effects of $\delta\approx0.15$--0.2 can therefore be neither confirmed nor excluded; they would need several hundred seeds. (v) The metal descriptors come from \texttt{mendeleev}, with zero electron affinity for elements with an unbound anion, which differs from the original Mordred values (Appendix~\ref{app:impl}).
